\documentclass[a4paper,12pt]{article}
\usepackage[fontset=none]{ctex}
\usepackage{geometry}
\usepackage{graphicx}
\usepackage{amsmath}
\usepackage{enumitem}
\usepackage{setspace}
\usepackage{fontspec}
\usepackage{booktabs}
\usepackage{caption}
\usepackage{array}
\usepackage{multirow}

% 字体设置（与 docx 模板一致）
\setCJKmainfont{Songti SC}          % 宋体 - 正文默认
\setCJKsansfont{Heiti SC}            % 黑体 - 一级标题
\setCJKmonofont[AutoFakeBold]{STFangsong}          % 仿宋 - 注意事项
\setmainfont{Times New Roman}        % 英文默认

% 页面设置
\geometry{left=2.5cm, right=2.5cm, top=2.5cm, bottom=2.5cm}

% 去掉页码
\pagestyle{empty}

\begin{document}

% ========== 封面 ==========
\begin{center}
    \vspace*{2cm}
    {\fontsize{36pt}{43pt}\selectfont \textbf{物理实验报告}}
    \vspace{3cm}
\end{center}

\begin{center}
    \fontsize{16pt}{24pt}\selectfont
    \textbf{实验名称：数字示波器的原理及应用\quad} \\[0.5em]
    \textbf{实验桌号：\_\_\_\_\_\_\_\_\_\_\_\_\_\_\_\_\_\_\_\_\_\_\_\_\_} \\[0.5em]
    \textbf{指导教师：\_\_\_\_\_\_\_\_\_\_\_\_\_\_\_\_\_\_\_\_\_\_\_\_\_}
\end{center}

\vspace{2cm}

\begin{center}
    \fontsize{14pt}{28pt}\selectfont
    \textbf{周次：\_\_\_\_\_\_\_\_\_\_\_\_\_\_\_\_\_\_\_\_\_\_\_\_\_\_\_\_\_\_\_\_\_} \\
    \textbf{姓名：\_\_\_\_\_\_\_\_\_\_\_\_\_\_\_\_\_\_\_\_\_\_\_\_\_\_\_\_\_\_\_\_\_} \\
    \textbf{学号：\_\_\_\_\_\_\_\_\_\_\_\_\_\_\_\_\_\_\_\_\_\_\_\_\_\_\_\_\_\_\_\_\_} \\[0.5em]
    \textbf{实验日期: \_\_\_\_年\_\_\_\_月\_\_\_\_日\quad 星期\_\_\_上/下午}
\end{center}

\vspace{1.5cm}

\begin{center}
    \fontsize{14pt}{20pt}\selectfont
    浙江大学物理实验教学中心
\end{center}

\newpage

% ========== 一、预习报告 ==========
{\CJKfontspec{Heiti SC}\fontsize{16pt}{24pt}\selectfont \textbf{一、预习报告}}

\vspace{0.5cm}
{\fontsize{14pt}{20pt}\selectfont \textbf{1. 实验综述}}

\vspace{0.3cm}

{\fontsize{12pt}{18pt}\selectfont

示波器是一种用途广泛的电子测量仪器，能观测电信号波形并测量其幅度、周期、频率和相位差等参量。本实验使用RIGOL DS1102型数字示波器和SG1020A信号发生器进行测量。

数字示波器的工作原理：待测信号经探头输入后，首先通过衰减器与放大器调节幅度至ADC的合适输入范围；然后由模数转换器（ADC）以一定采样率将模拟信号转换为离散的数字采样点并存入存储器；触发模块采用施密特触发方式，将触发信号转换为同频率方波以确定采样的时间起止点，实现与待测信号的同步，使波形稳定显示；最后由数据处理器对采样点进行内插、测量等处理后在液晶屏上显示波形。示波器的关键指标包括带宽、采样率和存储深度。

实验内容包括：（1）使用直读法和光标法测量正弦波、方波、三角波的峰峰值与频率；（2）利用X-Y模式观察李萨如图形，根据$f_y : f_x = N_x : N_y$的关系测量未知信号频率；（3）测量二极管正向导通电压及伏安特性；（4）使用驻波法测量超声波声速。

}

\vspace{0.5cm}

{\fontsize{14pt}{20pt}\selectfont \textbf{2. 实验重点}}

\vspace{0.3cm}

{\fontsize{12pt}{18pt}\selectfont

掌握数字示波器的基本结构和工作原理，熟悉垂直挡位、水平时基和触发电平的调节方法；学会使用直读法、光标法和自动测量功能测量信号的幅度、频率等参量；掌握X-Y模式下利用李萨如图形测量未知信号频率的方法。

}

\vspace{0.5cm}

{\fontsize{14pt}{20pt}\selectfont \textbf{3. 实验难点}}

\vspace{0.3cm}

{\fontsize{12pt}{18pt}\selectfont

正确选择触发源、耦合方式并调节触发电平以获得稳定波形；在李萨如图形观测中准确判别交点数$N_x$和$N_y$以确定频率比；理解带宽、采样率与存储深度三者之间的制约关系及其对测量精度的影响。

}

\newpage

% ========== 二、原始数据 ==========
{\CJKfontspec{Heiti SC}\fontsize{16pt}{24pt}\selectfont \textbf{二、原始数据}}

\vspace{0.3cm}
{\fontsize{12pt}{18pt}\selectfont （将有学生和老师签名的"自备数据记录草稿纸"的扫描或手机拍摄图粘贴在下方）}

\vspace{0.5cm}

% TODO: 插入原始数据纸照片
\begin{center}
\fbox{\parbox{0.85\textwidth}{\centering\vspace{6cm}{\fontsize{14pt}{20pt}\selectfont 【原始数据纸照片占位】}\vspace{6cm}}}
\end{center}

\newpage

% ========== 三、结果与分析 ==========
{\CJKfontspec{Heiti SC}\fontsize{16pt}{24pt}\selectfont \textbf{三、结果与分析}}

\vspace{0.5cm}

{\fontsize{14pt}{20pt}\selectfont \textbf{1. 数据处理与结果}}

\vspace{0.3cm}

{\fontsize{12pt}{18pt}\selectfont

% -------------------- [1] 比较法验证 fy = n*fx --------------------
\textbf{[1] 用比较法验证$f_y = nf_x$}

\vspace{0.2cm}

调节扫描时基因数为 \underline{\quad\quad} ms/DIV，因此$f_x$理论值 = \underline{\quad\quad} Hz。信号发生器输出正弦波，峰峰值 \underline{\quad\quad}\,V。

\vspace{0.3cm}

{\fontsize{10pt}{14pt}\selectfont
\begin{tabular}{c|c|c|c|c|c|c}
    \toprule
    波形个数 $n$ & 1 & 2 & 3 & 4 & 5 & 6 \\
    \midrule
    $f_y$ / Hz & & & & & & \\[0.5em]
    $f_{x\text{测定}} = f_y / n$ / Hz & & & & & & \\[0.5em]
    \bottomrule
\end{tabular}
}

\vspace{0.3cm}

$f_x$测定平均值：
\[
\bar{f}_{x\text{测定}} = \frac{1}{6}\sum_{i=1}^{6} f_{xi} = \underline{\quad\quad\quad} \text{ Hz}
\]

相对误差：
\[
E = \frac{|\bar{f}_{x\text{测定}} - f_{x\text{理论}}|}{f_{x\text{理论}}} \times 100\% = \underline{\quad\quad\quad}\%
\]

\vspace{0.5cm}

% -------------------- [2] 李萨如图形 --------------------
\textbf{[2] 用李萨如图形测量未知信号的频率}

\vspace{0.2cm}

理论值：$f_y$理论 = \underline{\quad\quad} Hz

\vspace{0.3cm}

{\fontsize{10pt}{14pt}\selectfont
\begin{tabular}{c|c|c|c|c|c}
    \toprule
    频率比 $f_y : f_x$ & $1:1$ & $1:2$ & $1:3$ & $2:1$ & $2:3$ \\
    \midrule
    垂直交点数 $N_y$ & & & & & \\[0.5em]
    水平交点数 $N_x$ & & & & & \\[0.5em]
    读出 $f_x$ / Hz & & & & & \\[0.5em]
    计算 $f_{y\text{测定}}$ / Hz & & & & & \\[0.5em]
    \bottomrule
\end{tabular}
}

\vspace{0.3cm}

频率比与交点数关系：
\begin{equation}
    f_y : f_x = N_x : N_y
\end{equation}

$f_y$测定平均值：
\[
\bar{f}_{y\text{测定}} = \frac{1}{5}\sum_{i=1}^{5} f_{yi} = \underline{\quad\quad\quad} \text{ Hz}
\]

相对误差：
\[
E = \frac{|\bar{f}_{y\text{测定}} - f_{y\text{理论}}|}{f_{y\text{理论}}} \times 100\% = \underline{\quad\quad\quad}\%
\]

\vspace{0.5cm}

% -------------------- [3] 二极管正向导通电压 --------------------
\textbf{[3] 二极管正向导通电压测量}

\vspace{0.2cm}

二极管型号：1N4007（硅二极管），分压电阻$R = 750$\,$\Omega$。信号发生器输出频率 \underline{\quad\quad}\,Hz、峰峰值 \underline{\quad\quad}\,V的正弦波信号，时基因数 \underline{\quad\quad}\,ms/DIV。

\vspace{0.2cm}

测量数据：

$V_{1\text{p-p}}$ = \underline{\quad\quad} V \qquad $V_{2\text{p}}$ = \underline{\quad\quad} V

\vspace{0.2cm}

二极管正向导通电压计算公式：
\begin{equation}
    V_{\text{导}} = \frac{V_{1\text{p-p}}}{2} - V_{2\text{p}}
\end{equation}

代入数据：
\[
V_{\text{导}} = \frac{\underline{\quad\quad}}{2} - \underline{\quad\quad} = \underline{\quad\quad} \text{ V}
\]

\vspace{0.5cm}

% -------------------- [4] 相位差测量 --------------------
\textbf{[4] RC电路相位差的测量}

\vspace{0.2cm}

电阻$R = 750$\,$\Omega$，电容$C = 0.47$\,$\mu$F。信号发生器输出频率 \underline{\quad\quad}\,Hz、峰峰值 \underline{\quad\quad}\,V 的正弦波信号，时基因数 \underline{\quad\quad}\,ms/DIV。

\vspace{0.2cm}

测量数据：

$T$ = \underline{\quad\quad} ms \qquad $\Delta t$ = \underline{\quad\quad} ms

\vspace{0.2cm}

相位差计算公式：
\begin{equation}
    \Delta\varphi = \frac{\Delta t}{T} \times 360°
\end{equation}

代入数据：
\[
\Delta\varphi = \frac{\underline{\quad\quad}}{\underline{\quad\quad}} \times 360° = \underline{\quad\quad\quad}
\]

理论值（$f = $ \underline{\quad\quad} Hz，$R = 750$ $\Omega$，$C = 0.47\,\mu$F）：
\begin{equation}
    \Delta\varphi_{\text{理论}} = \arctan(2\pi f R C) = \underline{\quad\quad\quad}
\end{equation}

}

\vspace{0.5cm}

{\fontsize{14pt}{20pt}\selectfont \textbf{2. 误差分析}}

\vspace{0.3cm}

{\fontsize{12pt}{18pt}\selectfont

\textbf{[1] 用比较法验证$f_y = nf_x$}

相对误差：
\[
E = \frac{|\bar{f}_{x\text{测定}} - f_{x\text{理论}}|}{f_{x\text{理论}}} \times 100\% = \underline{\quad\quad\quad}\%
\]

波形有一定的宽度，对于线条中间位置的判断不够准确，可能导致测量出现误差；直读法读取格数时存在人眼视差，且示波器显示器像素有限，读数存在量化误差。

\vspace{0.3cm}

\textbf{[2] 用李萨如图形测量未知信号的频率}

相对误差：
\[
E = \frac{|\bar{f}_{y\text{测定}} - f_{y\text{理论}}|}{f_{y\text{理论}}} \times 100\% = \underline{\quad\quad\quad}\%
\]

由于实验使用的是交流电信号源，李萨如图形无法保持完全静止，只能近似到尽可能缓慢的变化，读取信号发生器频率时存在微小偏差。此外，信号发生器频率的有限分辨率也会引入误差。

\vspace{0.3cm}

\textbf{[3] 二极管正向导通电压测量与相位差测量}

$V_{\text{导}}$应处于硅二极管正常正向导通电压范围（$0.5\sim0.7$\,V）附近。可能的误差来源：两个正弦波信号的振幅不相等，导致波峰等判断可能存在误差；光标法测定时，由于光标存在宽度，无法精准定位到相切位置。相位差测量中，$\Delta t$的读数精度受限于示波器的时间分辨率，也会引入一定误差。

}

\vspace{0.5cm}

{\fontsize{14pt}{20pt}\selectfont \textbf{3. 实验探讨}}

\vspace{0.3cm}

{\fontsize{12pt}{18pt}\selectfont

本次实验通过比较法验证频率关系、观察李萨如图形测量未知频率、测量二极管正向导通电压和RC电路相位差等内容，系统练习了数字示波器的使用方法，加深了对触发同步原理和信号测量方法的理解，为后续电学实验奠定了基础。

}

\newpage

% ========== 四、思考题 ==========
{\CJKfontspec{Heiti SC}\fontsize{16pt}{24pt}\selectfont \textbf{四、思考题}}

\vspace{0.5cm}

{\fontsize{12pt}{18pt}\selectfont

\textbf{1. 数字示波器为什么能显示被测信号的波形？}

\vspace{0.2cm}

数字示波器通过以下过程显示波形：（1）水平方向上，触发模块产生与待测信号同步的时钟信号，确定采样的时间起止点，将时间维度映射为屏幕水平方向的位移；（2）垂直方向上，ADC对待测信号进行高速采样，将连续的模拟电压转换为离散数字量，反映为屏幕垂直方向的像素位置；（3）数据处理器将存储器中的采样点经内插重建后输出至液晶屏逐像素点亮，由于刷新速率足够高，人眼即可看到连续稳定的波形图像。正确触发保证每次刷新的波形起始相位一致，从而使图像稳定不漂移。

\vspace{0.5cm}

\textbf{2. 在观察李萨如图形时，波形为什么会来回翻转，翻转的速度受什么因素影响？}

\vspace{0.2cm}

李萨如图形的形成要求两个正弦信号的频率比精确保持恒定且相位差固定。实际实验中，由于信号发生器的频率稳定性有限，两路信号的频率无法做到绝对恒定的整数比，存在微小的频率偏差$\Delta f$。这使得两信号的相位差随时间缓慢变化，导致李萨如图形不断翻转演化。

翻转速度主要受频率偏差$\Delta f$的影响：偏差越大，相位差变化越快，图形翻转越迅速；当$\Delta f \to 0$时，图形趋于稳定；当频率比恰好为简单整数比（$\Delta f = 0$）时，图形完全静止。此外，信号源的频率分辨率和短期稳定性也影响翻转快慢。

\vspace{0.5cm}

\textbf{3. 当波形左移或右移时，应如何调整触发使波形稳定？}

\vspace{0.2cm}

波形左移或右移说明扫描速率与被测信号的频率之间未达到同步。具体而言：波形左移表示被测信号的周期略小于扫描时基周期，波形右移则相反。

调整方法：（1）调节水平控制区的TIME/DIV旋钮（SCALE），改变扫描时基因数，使水平周期与信号周期匹配；（2）调节触发电平（LEVEL）旋钮，使触发电平处于信号幅度范围内的合适位置；（3）确认触发源选择正确（与输入信号通道一致）；（4）选择合适的触发斜率（上升沿或下降沿），使每次触发点位于信号的同一相位。当以上条件满足后，波形即可稳定显示。

}

\vfill

% ========== 注意事项 ==========
\noindent\rule{\textwidth}{0.4pt}

{\CJKfontspec[AutoFakeBold]{STFangsong}\fontsize{12pt}{18pt}\selectfont
\textbf{注意事项：}
\begin{enumerate}[label=\arabic*., leftmargin=2em, nosep]
    \item 用PDF格式上传"实验报告"，文件名：学生姓名+学号+实验名称+周次。
    \item "实验报告"必须递交在"学在浙大"的本课程的对应实验项目的"作业"模块内。
    \item "实验报告"成绩必须在"浙江大学物理实验教学中心网站"-"选课系统"内查询。
    \item 教学评价必须在"浙江大学物理实验教学中心网站"-"选课系统"内进行，学生必须进行教学评价，才能看到实验报告成绩，教学评价必须在本次实验结束后3天内进行。
\end{enumerate}
}

\vspace{0.5cm}
\begin{center}
    {\CJKfontspec[AutoFakeBold]{STFangsong}\fontsize{12pt}{18pt}\selectfont \textbf{浙江大学物理实验教学中心制}}
\end{center}

\end{document}
